\documentclass{article}
\usepackage[utf8]{inputenc}
\usepackage{amsmath,amssymb}
\usepackage[most]{tcolorbox}
\usepackage{geometry}
\geometry{margin=2.5cm}

\newcommand{\step}[1]{\par\noindent\hangindent=3.5em\hangafter=1%
  \makebox[3.5em][l]{\textbf{#1.}}}
\newcommand{\steptag}[1]{\hfill{\small\textsf{[#1]}}}

\newtcolorbox{proofbox}[1]{
  colback=gray!5, colframe=gray!60,
  fonttitle=\bfseries, title={#1},
  breakable, enhanced,
  left=4pt, right=4pt, top=4pt, bottom=4pt,
  before skip=10pt, after skip=10pt
}

\begin{document}

\begin{proofbox}{Route F scalar-header positivity: there exists one global...}
\step{1} Route F scalar-header positivity: there exists one global witness package W\_RF supplied by lem-routef-raw-factor-setting-formation such that, writing K for its scalar (1.6), rho\_fac for its scalar (1.7), and eta\_K := min\{rho\_fac, (24*K)\textasciicircum{}(-1), 1\} for its scalar (1.8), K is finite with K >= 1, rho\_fac > 0, and eta\_K > 0, these scalars are universal and independent of H, Phi, eta, n, amplification level, simple-block count, and block dimensions, and eta\_K <= rho\_fac <= rho\_2 <= rho\_T <= rho\_id\textasciicircum{}corr, rho\_2 <= rho\_Delta', rho\_2 <= rho\_Delta, rho\_fac <= rho\_DeltaUpsilon <= rho\_Upsilon <= rho\_Upsilon', rho\_fac <= rho\_mult, and rho\_fac <= rho\_UpsilonDelta. \steptag{validated} \\[6pt]
\step{1.1} The external lem-routef-raw-factor-setting-formation supplies a scalar header W\_RF before all input quantifiers, independent of H, Phi, eta, dimension, amplification level, and block data, with eta\_A>0, C\_E finite, epsilon\_E>0, C\_theta=12*(sqrt(2)-1), C\_A=20+(211/8)*C\_theta, rho\_theta=1/8, rho\_AI=eta\_A, and every remaining scalar defined by def-routef-raw-factor-setting equations (1.1)-(1.8). \steptag{validated} \\[4pt]
\step{1.2} For every real scalar header W\_RF having the primitive facts and definitions listed in node 1.1, equations (1.1)-(1.8) imply that K is finite with K>=1, rho\_fac>0, and eta\_K>0. \steptag{validated} \\[6pt]
\step{1.2.1} The primitive layer is positive and finite: sqrt(2)>1 makes C\_theta=12*(sqrt(2)-1)>0 and finite; hence C\_A=20+(211/8)*C\_theta>0 and finite. Since C\_E is a finite real scalar, bar\_C\_E=max\{1,C\_E\} is finite and at least 1, so C\_V=bar\_C\_E*C\_A and C\_T=C\_theta+3*C\_V are positive and finite. Also rho\_theta=1/8>0, rho\_AI=eta\_A>0, and epsilon\_E/C\_A>0; the two reciprocal entries in (1.1) have positive finite denominators. Thus every entry defining rho\_T is positive, so rho\_T>0. \steptag{validated} \\[4pt]
\step{1.2.2} Sequentially applying (1.2) to the positive finite quantities from node 1.2.1 gives rho\_unit=rho\_T>0, rho\_id>0, rho\_id\textasciicircum{}corr>0, rho\_prod=rho\_T>0, C\_Delta'>0 finite, rho\_Delta'>0, C\_Delta>0 finite, rho\_Delta>0, C\_2>0 finite, rho\_2>0, rho\_DeltaPhi>0, C\_3>0 finite, and rho\_3>0: each new radius is a finite minimum of already-positive terms, and the only new reciprocal has denominator 2*(C\_T+C\_Delta')>0. \steptag{validated} \\[4pt]
\step{1.2.3} Equations (1.3)-(1.5), evaluated after node 1.2.2, give positive finite C\_N,C\_R,C\_L,C\_Upsilon', then rho\_Upsilon'>0 because its six entries are positive (including (2*C\_R)\textasciicircum{}(-1)); next C\_Upsilon>0 finite and rho\_Upsilon>0 because its entries are positive (including [2*(C\_T+C\_Upsilon')]\textasciicircum{}(-1)); finally rho\_DeltaUpsilon, rho\_mult, and rho\_UpsilonDelta are positive as minima of positive radii. \steptag{validated} \\[4pt]
\step{1.2.4} By (1.6), K is the maximum of finitely many finite real numbers and includes the entry 1, so K is finite and K>=1. Validated node 1.2.2 supplies rho\_2>0, while validated node 1.2.3 supplies rho\_DeltaUpsilon>0, rho\_mult>0, and rho\_UpsilonDelta>0. These are exactly the four entries in (1.7), hence rho\_fac>0. Since K>=1, (24*K)\textasciicircum{}(-1)>0; therefore (1.8) makes eta\_K=min\{rho\_fac,(24*K)\textasciicircum{}(-1),1\}>0. \steptag{validated} \\[6pt]
\step{1.2.4.1} By validated nodes 1.2.1, 1.2.2, and 1.2.3, the constituent scalars C\_theta, C\_Delta, C\_2, and C\_Upsilon are finite real numbers. Hence the three nonconstant expressions in (1.6), namely C\_theta+C\_Delta+2*C\_Upsilon, C\_Upsilon+2*(C\_2+C\_theta+C\_Delta), and C\_Upsilon+2*C\_Delta, are finite real numbers. Therefore (1.6) defines K as the maximum of four finite real numbers: 1 and those three expressions. Thus K is finite; since 1 is one of the four entries, K>=1. \steptag{validated} \\[4pt]
\step{1.2.4.2} Validated node 1.2.2 gives rho\_2>0, and validated node 1.2.3 gives rho\_DeltaUpsilon>0, rho\_mult>0, and rho\_UpsilonDelta>0. These are exactly the four entries in (1.7), so their finite minimum rho\_fac is strictly positive. \steptag{archived} \\[4pt]
\step{1.2.4.3} From child 1.2.4.1, K>=1, so 24*K>=24>0 and therefore the real reciprocal (24*K)\textasciicircum{}(-1)>0. From child 1.2.4.2, rho\_fac>0, while 1>0. Thus all three entries of the minimum in (1.8) are strictly positive, and consequently eta\_K=min\{rho\_fac,(24*K)\textasciicircum{}(-1),1\}>0. Together with 1.2.4.1 and 1.2.4.2 this proves the statement of node 1.2.4. \steptag{archived} \\[4pt]
\step{1.2.4.4} Validated node 1.2.2 gives rho\_2>0, and validated node 1.2.3 gives rho\_DeltaUpsilon>0, rho\_mult>0, and rho\_UpsilonDelta>0. These are exactly the four entries of (1.7), so their finite minimum rho\_fac is strictly positive. By child 1.2.4.1, K>=1; hence 24*K>=24>0 and the real reciprocal (24*K)\textasciicircum{}(-1)>0. Since also 1>0, every entry of the minimum in (1.8) is strictly positive, so eta\_K=min\{rho\_fac,(24*K)\textasciicircum{}(-1),1\}>0. Together with child 1.2.4.1 this proves every conclusion of node 1.2.4. \steptag{validated} \\[4pt]
\step{1.3} For every scalar header governed by def-routef-raw-factor-setting equations (1.1)-(1.8), the coordinate inequalities for the displayed finite minima imply eta\_K<=rho\_fac<=rho\_2<=rho\_T<=rho\_id\textasciicircum{}corr, rho\_2<=rho\_Delta', rho\_2<=rho\_Delta, rho\_fac<=rho\_DeltaUpsilon<=rho\_Upsilon<=rho\_Upsilon', rho\_fac<=rho\_mult, and rho\_fac<=rho\_UpsilonDelta. \steptag{validated} \\[6pt]
\step{1.3.1} Equation (1.8) gives eta\_K<=rho\_fac, and (1.7) gives rho\_fac<=rho\_2. Equation (1.2) gives rho\_2<=rho\_prod=rho\_T, rho\_2<=rho\_Delta', and rho\_2<=rho\_Delta. Finally rho\_T is the minimum in (1.1) of the three entries rho\_theta, rho\_AI, epsilon\_E/C\_A together with two more entries, whereas rho\_id\textasciicircum{}corr=min\{rho\_theta,rho\_AI,epsilon\_E/C\_A\}; hence rho\_T<=each of the three entries and therefore rho\_T<=rho\_id\textasciicircum{}corr. \steptag{validated} \\[4pt]
\step{1.3.2} Equation (1.7) gives rho\_fac<=rho\_DeltaUpsilon. Equation (1.5) gives rho\_DeltaUpsilon<=rho\_Upsilon. The definition of rho\_Upsilon in (1.5) includes rho\_Upsilon' among its minimum entries, so rho\_Upsilon<=rho\_Upsilon'. Thus rho\_fac<=rho\_DeltaUpsilon<=rho\_Upsilon<=rho\_Upsilon'. \steptag{validated} \\[4pt]
\step{1.3.3} The two remaining coordinate inequalities follow directly from (1.7), whose minimum entries include rho\_mult and rho\_UpsilonDelta: rho\_fac<=rho\_mult and rho\_fac<=rho\_UpsilonDelta. \steptag{validated} \\[4pt]
\step{1.4} Because W\_RF is chosen independently before every input quantifier and all K, rho\_fac, and eta\_K are obtained from its scalar fields alone by the fixed formulas (1.1)-(1.8), these three scalars are universal and independent of H, Phi, eta, n, amplification level, simple-block count, and block dimensions. \steptag{validated} \\[4pt]
\end{proofbox}

\end{document}
